\subsection{Related Works}

Our work builds on the intersection of strategic machine learning, sender-receiver games, and voluntary disclosure in predictive modeling. Early studies in strategic classification, such as \cite{hardt2016strategic}, formalized settings where agents manipulate features to influence outcomes, showing that robust classifiers can mitigate gaming in linear models. This was extended in \cite{ghalme2021strategic} to dark (unknown) environments, where agents' responses are not fully observable, leading to equilibrium analyses under uncertainty.

Sender-receiver games provide a foundational framework for our analysis. \cite{blume2002learning} examined learning dynamics in these games using experimental data, highlighting how repeated interactions evolve into efficient signaling. Dynamic variants, like those in \cite{renault2013dynamic}, incorporate Markov chains for state evolution, akin to our noisy channel model. Evolutionary perspectives, as in \cite{skyrms2009evolution}, explore multi-agent signaling, while \cite{blume1998learning} compared stimulus-response and belief-based learning, informing our online regret bounds.

Voluntary disclosure literature addresses selective revelation. \cite{aghamolla2019voluntary} modeled evolving news in dynamic settings, deriving conditions for disclosure thresholds. \cite{mcdonough2023voluntary} analyzed investor disclosures, predicting good-news bias, which parallels our heterogeneous preferences. Empirical studies, such as \cite{lu2024determinants}, used XGBoost and XAI for predicting disclosure patterns, emphasizing transparency in strategic contexts.

Stackelberg equilibria are central to our game-theoretic approach. \cite{jin2022incentive} characterized equilibria in strategic classification and regression, designing incentives to prevent cheating. \cite{sundaram2023pac} generalized PAC-learning for heterogeneous agents, bridging classification and regression. In adversarial prediction, \cite{bruckner2011stackelberg} framed problems as Stackelberg games, solving for optimal decision rules.

Online learning with strategic agents extends our dynamic framework. \cite{yin2022online} studied allocation under strategic behavior, deriving regret-minimizing algorithms. \cite{wu2025learning} proposed online principal-agent models for unknown types, achieving sublinear regret. Decentralized settings in \cite{cohen2025decentralized} model coalition formation among selfish agents, aligning with our adaptive disclosure policy.

While these works advance strategic aspects in classification \cite{trachtenberg2025strategic} and regression \cite{shavit2020causal}, they often assume known environments or linear models. Our contributions provide formal sufficiency and necessity conditions under noisy channels and heterogeneous preferences, with theoretical bounds for online adaptation, filling a gap in generalized predictive settings.
